\section{The audit prompted by the PLoP paper}\label{sec:case-vetting}

%% Provenance for the numbers in this section, which are NOT paper macros:
%% every figure is quoted from a dated internal record rather than regenerated
%% by a script in empirics-futon/ -- a departure from WR-8, flagged as in
%% sec-boundary. The records:
%%   futon2/holes/labs/wm-contract/NOTE-cascade-consumers-census.md   (07-05/06/09 dates)
%%   futon2/holes/labs/wm-contract/claude-15-repl-buffer-snapshot-2026-08-31.txt
%%   vetting/OBLIGATIONS-REVERIFY-2026-08-31.md                        (24 obligations)
%%   futon2/holes/problems/P-defect-classes.md                         (9 classes)
%%   futon2/holes/problems/BUILD-ledger.md, WM-RUN2 gate 2026-08-30    (3 of 9 hops)
%%   futon3 library harvest be8c707 review ad0b004                     (spider counts)
%% These are frozen historical facts about one episode, not live joins; freezing
%% them by hand and naming the record is the honest form here.


The PLoP catalogue was treated as finished on 2026-08-21. On 2026-08-30,
asking who consumed the output of its Learn stage reopened the implementation
question. The census found that Learn's input had stopped arriving on July~6,
the nightly observer had last run on July~5, and writeback deposits had stopped
that week. The operator's question detected the failure; no standing monitor
had reported the producer and consumer falling silent.

The following day's vetting checked a twenty-four-row obligation ledger against
published PDFs and found two unsupported closures. The rows were subsequently
resolved by corrections, new evidence or explicit withdrawals. That was a
repair of the paper's obligations, not proof that the machine worked. An
instrumented run separately agreed with the described wiring on only three of
nine hops. These dated findings began the repair programme now represented by
the issue board.

\subsection*{What the audit checks}

The investigation separates five questions. \emph{Nouns}: do the records denote
the objects the claim requires? \emph{Verbs}: do producers and consumers perform
the claimed operations? \emph{Organisation}: do their handoffs compose, with
compatible types and a responsible owner? \emph{Evidence}: can a check reject
a false claim, and does its receipt refer to the artifact being discussed?
Finally, who is present to act on a refusal or supply the required judgement?
A missing answer is an obligation, not permission to substitute a nearby one.

This decomposition explained how passing checks could coexist with a broken
loop. A theorem about an object does not establish that a live caller supplies
that object. A generated artifact is not evidence that its consumer runs.
Likewise, resolving a documentation defect does not resolve the implementation
defect it describes. Independent review must check the asserted behaviour and
its evidence, not merely the summary's status.

The record also made the investigation resumable when its original agent
session became unavailable. Another agent could continue from dated ledgers
and problem records. This supports a practical requirement: an unfinished task
must retain its current evidence, next step and owner outside the conversation
that happened to discover it.

\subsection*{Apparatus patterns and their limits}\label{sec:apparatus-quality}

The new \texttt{futon3/library/apparatus/} collection describes mechanisms for
keeping the investigative machinery answerable to its work. Three are
especially relevant here. \emph{Done is observed running} requires evidence
that a change is loaded, connected and producing its intended effect.
\emph{Monitors measure the work} checks durable progress of the actual task,
not the coordinator's heartbeat, and calls for commissioning against both a
stall and healthy progress. \emph{Evidence to disposition once} requires a
shared interpretation of raw observations before different consumers derive
their own decisions.

Other entries address authority, pin updates, deadlines and repair citations.
Their context matters: the collection's field trial is the APM apparatus repair
programme, documented in
\path{futon3c/holes/technotes/TN-apm-principle-scoring-ledger.md}.
That ledger records principles that helped, misled their users or needed
revision. It is evidence about those episodes, not a guarantee that every
pattern is installed in the War Machine. A transfer claim needs its own
consumer, failure control and witnessed result. In particular, keeping a pin
current must include the required independent review; it does not authorise
an implementer to approve their own evidence.

These patterns are candidates for a focused addition to the PLoP paper. Their
value is to explain concrete failures and repairs, rather than add a second
catalogue of desirable practices. The paper's existing component witnesses
remain useful, but the larger claim awaits an integrated run at an explicit
contract and source basis.

\subsection*{From audit to action}

The board below is the current work view. It preserves disagreements among
sources instead of resolving them by presentation, distinguishes unowned work
from an absent carrying packet, and displays evidence age. A readiness result
admits an attempt; a run-gated obligation needs that attempt's evidence.
Neither a green readiness line nor a completed repair count is a completed,
compliant run. The narrative explains the audit; the board owns its changing
status.

\clearpage
\input{sec-issue-board-generated}
\clearpage
